% Generated draft: requires mathematical and source review.
\documentclass[11pt]{article}
\usepackage[margin=25mm]{geometry}
\usepackage{fontspec}
\setmainfont{Latin Modern Roman}
\usepackage{amsmath,amssymb,mathtools,unicode-math}
\setmathfont{Latin Modern Math}
\usepackage{enumitem,needspace,xurl,hyperref}
\setlength{\parindent}{0pt}
\setlength{\parskip}{5pt}
\setlength{\emergencystretch}{3em}
\newcommand{\N}{\mathbb{N}}
\newcommand{\Z}{\mathbb{Z}}
\newcommand{\Q}{\mathbb{Q}}
\newcommand{\R}{\mathbb{R}}
\newcommand{\C}{\mathbb{C}}
\newcommand{\F}{\mathbb{F}}
\newcommand{\A}{\mathbb{A}}
\newcommand{\PP}{\mathbb{P}}
\newcommand{\E}{\mathbb{E}}
\newcommand{\eps}{\varepsilon}
\newcommand{\dd}{\,\mathrm{d}}
\begin{document}
\section*{UnsolvedMath ID: 30005755}
\textbf{Source problem code:} catalogue entry
\subsection{Definitions and mathematical statement}
Let \(k\) be an algebraically closed field and let \(A\) be a finite-dimensional \(k\)-algebra, as required by the supplied original passage. Write \(\operatorname{proj}A\) for the category of finitely generated projective right \(A\)-modules and put
\[
V_A=K_0(\operatorname{proj}A)\otimes_{\Z}\R.
\]
The use of right modules is an explicit convention for this draft; the supplied extract does not specify handedness, which must be checked against the references.

For \(\eta\in V_A\) and a finite-dimensional right \(A\)-module \(M\), define \(\eta(M)\) by real-linear extension of
\[
[P](M)=\dim_k\operatorname{Hom}_A(P,M).
\]
Define the two weak numerical torsion classes by
\[
\begin{aligned}
\overline{\mathcal T}_{\eta}
 &=\{M:\eta(Q)\geq 0\text{ for every quotient module }Q\text{ of }M\},\\
\overline{\mathcal F}_{\eta}
 &=\{M:\eta(L)\leq 0\text{ for every submodule }L\subseteq M\}.
\end{aligned}
\]
Two weights are TF-equivalent when both of these classes agree. Thus
\[
[\theta]_{\mathrm{TF}}
 =\{\eta\in V_A:
 \overline{\mathcal T}_{\eta}=\overline{\mathcal T}_{\theta}
 \text{ and }
 \overline{\mathcal F}_{\eta}=\overline{\mathcal F}_{\theta}\}.
\]

For an integral class \(\gamma\in K_0(\operatorname{proj}A)\), choose projectives \(P_\gamma^-\) and \(P_\gamma^+\) with no common indecomposable summands such that
\[
\gamma=[P_\gamma^+]-[P_\gamma^-].
\]
The generic canonical decomposition of \(\gamma\) is the decomposition into generically indecomposable projective-presentation classes obtained from a Zariski-generic map in \(\operatorname{Hom}_A(P_\gamma^-,P_\gamma^+)\). Write \(\operatorname{ind}(\gamma)\) for the set of its summand classes, ignoring multiplicities, and set
\[
\operatorname{ind}(\N\theta)
 :=\bigcup_{n\geq 1}\operatorname{ind}(n\theta).
\]
Here \(\operatorname{ind}(0)=\varnothing\); including \(n=0\) consequently has no effect. These are the standard projective-presentation interpretations of notation left undefined in the supplied extract and require source confirmation.

For \(S\subseteq V_A\), let
\[
\operatorname{cone}(S)
 =\left\{\sum_{i=1}^{r}a_i s_i:
 r\geq 0,\ a_i\in\R_{\geq0},\ s_i\in S\right\}.
\]
The empty sum is zero. Let \(\operatorname{relint}\) denote interior relative to the real linear span of this cone; no closure is taken before forming the relative interior.

The requested conjecture, identified as Conjecture 2 in the supplied passage, is: for every algebraically closed field \(k\), every finite-dimensional \(k\)-algebra \(A\), and every \(\theta\in K_0(\operatorname{proj}A)\), does one have
\[
[\theta]_{\mathrm{TF}}
 =\operatorname{relint}\!\left(
 \operatorname{cone}\bigl(\operatorname{ind}(\N\theta)\bigr)
 \right)?
\]
This includes \(\theta=0\), with \(\operatorname{cone}(\varnothing)=\{0\}\) and \(\operatorname{relint}\{0\}=\{0\}\).

The supplied passage attributes to its Theorem 1 the inclusion of the canonical-decomposition cone interior in \([\theta]_{\mathrm{TF}}\), and also asserts the inclusion for \(\operatorname{ind}(\N\theta)\). It states that \(\operatorname{cone}(\operatorname{ind}\theta)\subseteq\operatorname{cone}(\operatorname{ind}(\N\theta))\) need not be equality, referring to its Theorem 5. The extracted question interprets the report's notation \(\operatorname{cone}^{\circ}\) as relative interior; its agreement with the report's exact convention, especially for the potentially infinite generating set, remains to be checked.
\subsection{Short English statement}
For a finite-dimensional algebra over an algebraically closed field, ask whether every real weight TF-equivalent to an integral weight \(\theta\) lies precisely in the relative interior of the cone generated by the generic canonical summands of all positive multiples of \(\theta\). The supplied curation labels the problem partially solved, with hereditary and \(E\)-tame cases proved; this is a dated supplied assessment, not an independently verified conclusion.
\subsection{Sources}
\begin{itemize}
\item Supplied local UnsolvedMath registry record, permanent ID \(30005755\), title ``Canonical Decomposition Cones and TF Equivalence,'' exact source problem code \texttt{catalogue entry}. The record attributes the problem to the Oberwolfach Reports Open Problems collection and gives proposed year 2024; no individual proposer is supplied.
\item \emph{Cluster Algebras and Its Applications} (2024), Oberwolfach Reports, pp.\ 69--136, DOI \texttt{10.4171/owr/2024/2}: \url{https://doi.org/10.4171/owr/2024/2}. This is both the supplied source report and the literature-source link labelled ``original OWR report.'' The supplied extract identifies the target as Conjecture 2 and includes context attributed to Theorems 1 and 5; it supplies no precise page for that passage.
\item Supplied literature source labelled ``published primary paper'': \url{https://doi.org/10.1112/plms.12639}. No paper title, authors, or theorem locator is supplied.
\item Literature assessment checked at \texttt{2026-08-22}: the supplied curation states that the equality is proved for hereditary and \(E\)-tame algebras, while the general finite-dimensional-algebra case remains unresolved. This is supplied dated curation, not independently verified here. The background explicitly requests independent verification of time-sensitive status and cited claims.
\item The registry labels the statement \texttt{corrected\_verified} and reports that the official report was checked directly. This records the supplied verification claim; it does not claim that the present draft independently visited or verified either reference.
\end{itemize}
\end{document}
